% Part: second-order-logic
% Chapter: syntax-and-semantics
% Section: size-of-domain

\documentclass[../../../include/open-logic-section]{subfiles}

\begin{document}

\olfileid[fr]{sol}{syn}{siz}

\olsection{Décrire les domaines infinis et au plus
  dénombrables}

Un ensemble~$M$ est \emph{infini au sens de
Dedekind} si et seulement s'il existe une
fonction !!{injective}~$f\colon M\to M$ qui
n'est pas !!{surjective}, c'est-à-dire telle
que $\ran{f}\neq M$. Dans le cadre usuel où
l'on admet l'axiome du choix dénombrable,
cette propriété équivaut à l'infinitude.
En logique du premier ordre, nous pouvons
considérer un !!{function}
unaire~$f$ et affirmer que la fonction
$\Assign{f}{M}$ qui lui est attribuée dans
une !!{structure}~$\Struct{M}$ est
!!{injective} et que
$\ran{\Assign{f}{M}}\neq\Domain{M}$ :
\[
\lforall[x][\lforall[y][(\eq[f(x)][f(y)]
\to\eq[x][y])]]\land
\lexists[y][\lforall[x][\eq/[y][f(x)]]].
\]
Si $\Struct{M}$ satisfait cet énoncé,
$\Assign{f}{M}\colon\Domain{M}\to\Domain{M}$
est !!{injective} sans être !!{surjective};
le !!{domain} est donc infini.
Réciproquement, dans le cadre précisé ci-dessus,
si $\Domain{M}$ est infini, une telle fonction
existe. Nous pouvons alors prendre cette fonction
pour $\Assign{f}{M}$, et $\Struct{M}$ satisfera
l'énoncé. Il faut toutefois que notre langage
contienne le symbole non logique~$f$ utilisé
à cette fin. En logique du second ordre, nous
pouvons simplement affirmer qu'une telle
fonction \emph{existe}. Le symbole~$f$ n'est
plus nécessaire; nous obtenons l'énoncé suivant
de la logique pure du second ordre :
\[
\fn{Inf}\ident
\lexists[u][(\lforall[x][\lforall[y]
[(\eq[u(x)][u(y)]\lif\eq[x][y])]]
\land\lexists[y][\lforall[x][\eq/[y][u(x)]]])].
\]
$\Sat{M}{\fn{Inf}}$ si et seulement si
$\Domain{M}$ est infini au sens de Dedekind;
dans le cadre usuel précisé ci-dessus, si et
seulement s'il est infini. Posons ensuite
$\fn{Fin}\ident\lnot\fn{Inf}$; dans ce
même cadre, $\Sat{M}{\fn{Fin}}$ si et seulement
si $\Domain{M}$ est fini. Aucun énoncé unique
de la logique pure du premier ordre ne peut
exprimer que le !!{domain} est infini,
quoiqu'un ensemble infini de tels énoncés
le puisse. Aucun ensemble d'énoncés de la
logique pure du premier ordre n'est satisfait
exactement dans les !!{structure}s de
!!{domain} fini.

\begin{prop}
En présence de l'axiome du choix dénombrable,
$\Sat{M}{\fn{Inf}}$ si et seulement si
$\Domain{M}$ est infini.
\end{prop}

\begin{proof}
$\Sat{M}{\fn{Inf}}$ si et seulement si,
pour une certaine assignation~$s$,
\[
\Sat{M}{\lforall[x][\lforall[y]
[(\eq[u(x)][u(y)]\lif\eq[x][y])]]
\land\lexists[y][\lforall[x][\eq/[y][u(x)]]]}[s].
\]
Dans ce cas, $s(u)$ est une fonction
!!{injective}, et il existe un élément
$y\in\Domain{M}$ qui n'appartient pas à son
image. Réciproquement, s'il existe une fonction
!!{injective}
$f\colon\Domain{M}\to\Domain{M}$
avec $\ran{f}\neq\Domain{M}$, nous pouvons
choisir l'assignation~$s$ de sorte que
$s(u)=f$. Cela établit directement l'équivalence
avec l'infinitude au sens de Dedekind; l'axiome
du choix dénombrable donne celle avec
l'infinitude ordinaire.
\end{proof}

Un ensemble non vide~$M$ est !!{enumerable}
s'il existe une énumération de ses !!{element}s
\[
m_0,m_1,m_2,\dots
\]
sans répétition, éventuellement finie.
Une telle énumération existe si et seulement
s'il existe un !!{element}~$z\in M$ et une
fonction $f\colon M\to M$ tels que
$z$, $f(z)$, $f(f(z))$, \dots, parcourent tous
les !!{element}s de~$M$.
Dans un sens, si l'énumération existe,
prenons $z=m_0$ et $f(m_k)=m_{k+1}$,
en posant $f(m_k)=m_k$ si $m_k$ est le dernier
!!{element}. Dans l'autre sens, l'orbite
$z$, $f(z)$, $f(f(z))$, \dots, couvre~$M$;
en supprimant les répétitions éventuelles,
elle en fournit une énumération.
Nous pouvons exprimer l'existence de~$z$ et
de~$f$ en logique du second ordre par un
énoncé vrai dans une !!{structure} si et
seulement si son !!{domain} est
!!{enumerable} :
\[
\fn{Count}\ident
\lexists[z][\lexists[u][\lforall[X]
[((X(z)\land
\lforall[x][(X(x)\lif X(u(x)))])
\lif\lforall[x][X(x)])]]].
\]

\begin{prop}
$\Sat{M}{\fn{Count}}$ si et seulement si
$\Domain{M}$ est !!{enumerable}.
\end{prop}

\begin{proof}
Supposons $\Domain{M}$ !!{enumerable}, et
soit $m_0,m_1,\dots$ une énumération sans
répétition. Définissons
$f(m_k)=m_{k+1}$, en fixant le dernier
!!{element} si l'énumération est finie,
et prenons $s(z)=m_0$, $s(u)=f$.
Montrons que
\[
\Sat{M}{\lforall[X][((X(z)\land
\lforall[x][(X(x)\lif X(u(x)))])
\lif\lforall[x][X(x)])]}[s].
\]
Soit $S\subseteq\Domain{M}$ quelconque et
supposons
$\Sat{M}{(X(z)\land\lforall[x]
[(X(x)\lif X(u(x)))])}
[\Subst{s}{S}{X}]$.
Alors $m_0=s(z)\in S$, et si $a\in S$,
on a aussi $f(a)=s(u)(a)\in S$.
Par récurrence,
$m_0,m_1=f(m_0),m_2=f(f(m_0)),\dots$
appartiennent tous à~$S$. Puisqu'ils épuisent
$\Domain{M}$, on a $S=\Domain{M}$, donc
$\Sat{M}{\lforall[x][X(x)]}
[\Subst{s}{S}{X}]$. Comme~$S$ était
arbitraire, $\Sat{M}{\fn{Count}}$.

Réciproquement, supposons
$\Sat{M}{\fn{Count}}$, c'est-à-dire que,
pour une certaine assignation~$s$,
\[
\Sat{M}{\lforall[X][((X(z)\land
\lforall[x][(X(x)\lif X(u(x)))])
\lif\lforall[x][X(x)])]}[s].
\]
Posons $m=s(z)$, $f=s(u)$ et
$S=\{m,f(m),f(f(m)),\dots\}$.
Par construction, $S$ est !!{enumerable}.
L'hypothèse donne
\[
\Sat{M}{(X(z)\land\lforall[x]
[(X(x)\lif X(u(x)))])
\lif\lforall[x][X(x)]}
[\Subst{s}{S}{X}].
\]
Or $\Sat{M}{X(z)}[\Subst{s}{S}{X}]$,
car $m=s(z)\in S$; et
$\Sat{M}{\lforall[x][(X(x)\lif X(u(x)))]}
[\Subst{s}{S}{X}]$, car $S$ est stable
par~$f$. L'antécédent de l'implication est
donc satisfait et, par l'hypothèse, son
conséquent l'est également :
$\Sat{M}{\lforall[x][X(x)]}
[\Subst{s}{S}{X}]$. Il s'ensuit que
$S=\Domain{M}$. Le !!{domain} est donc
!!{enumerable}, puisque $S$ l'est.
\end{proof}

\begin{prob}
L'énoncé $\fn{Inf}\land\fn{Count}$ est vrai
exactement dans les !!{domain}s
!!{denumerable}s. Modifiez la définition de
$\fn{Count}$ pour obtenir un autre énoncé
qui exprime directement que le !!{domain} est
!!{denumerable}, puis démontrez-le.
\end{prob}

\begin{editorial}
La source assimile l'infinitude au sens de Dedekind
à l'infinitude sans préciser l'hypothèse de choix
nécessaire à cette réciproque. Nous indiquons
l'axiome du choix dénombrable, conformément au
cadre dans lequel ces notions coïncident.
L'énumération par une orbite demande aussi de
supprimer les répétitions éventuelles, et le
dernier élément d'une énumération finie doit
rester fixé par~$f$ dans la première moitié de
la preuve. Enfin, nous notons~$S$ les
sous-ensembles du !!{domain} que la source
appelle~$M$, afin de les distinguer de la
!!{structure}~$\Struct{M}$.
\end{editorial}

\end{document}
